\scanpage{87}\mspage{T-1}

\begin{center}
ELLIPTIC CURVES - FORMULARY

(Free translation of an unpublished manuscript by TATE)
\end{center}

\figure{\begin{tikzpicture}[x=0.016cm,y=0.016cm]
\draw[gray,thin] (0,0) .. controls (8,0) and (13,9) .. (22,14) .. controls (35,29) and (56,42) .. (77,43) .. controls (93,43) and (95,46) .. (97,50) -- (101,46);
\draw[gray,thin] (25,15) .. controls (32,0) and (43,-13) .. (47,-27);
\end{tikzpicture}}

1. \uline{General formulas}

If one starts with an elliptic c\deleted{\unsure{a|u}}urve $E$, equipped with an origin $O$, and a form of the first kind $\pi\ne 0$, one can find $x$ and $y$ on $E$, such that $x$ (resp. $y$) has a double (resp. triple) pole at $O$, and is regular elsewhere; moreover, if $\pi=dt+\ldots$, one can arrange (multiplying $x$ and $y$ by constants if necessary) that $x\handwritten{=}{t}^{-2}+\ldots$, $y=-t^{-3}+\ldots$

In the projective embedding defined by the divisor $3(O)$, the equation of the curve $E$ is written:
{\makeatletter\tagsleft@true\makeatother
\[
 y^{2}+a_{1}xy+a_{3}y=x^{3}+a_{2}x^{2}+a_{4}x+a_{6}
 \tag*{$(\ast)$}
\]
}

(Homogeneity: $y$ has weight $3$, $x$ weight $2$, and the $a_{i}$ weight $i$. This means that, if one replaces $\pi$ by $u\pi$, $x$ is replaced by $u^{-2}x$, $y$ by $u^{-3}y$, and so on.)

Conversely, let us start with an equation of the form $(\ast)$. We shall set:
\[
\begin{aligned}
& b_{2}=a_{1}^{2}+4a_{2} \qquad b_{4}=a_{1}a_{3}+2a_{4} \qquad b_{6}=a_{3}^{2}+4a_{6} \\
& -b_{8}=a_{1}a_{3}a_{4}+a_{4}^{2}-a_{1}^{2}a_{6}-a_{2}a_{3}^{2}-4a_{2}a_{6} \\
& b_{8}=b_{2}a_{6}+a_{2}a_{3}^{2}-a_{1}a_{3}a_{4}-a_{4}^{2}=a_{2}b_{6}+a_{1}^{2}a_{6}-a_{4}^{2}-a_{1}a_{3}a_{4} \\
& \hspace{4em}-4\,b_{8}=b_{4}^{2}-b_{2}b_{6}
\end{aligned}
\]
% archive sheet number: 86
